\documentclass[a4paper, 12pt, table]{article}
\usepackage[margin=0.8in]{geometry}
\usepackage{amsmath}
\numberwithin{equation}{section}
\usepackage{amsfonts}
\usepackage{xcolor}
\usepackage{ amssymb }
\usepackage{ dsfont }
\usepackage{caption}
\usepackage{graphicx}
\usepackage{pgfplots}
\usepackage{ upgreek }
\usepackage{enumitem}
\usepackage{hyperref}
\usepackage{amsthm}
\usepackage[english]{babel}
\usepackage{hyperref}
\usepackage{apacite}
\newtheorem{theorem}{Theorem}[section]\newtheorem{corollary}{Corollary}[theorem]\newtheorem{lemma}[theorem]{Lemma}\theoremstyle{definition}\newtheorem{definition}{Definition}[section]\theoremstyle{remark}\newtheorem*{remark}{Remark}

\title{Statistical Profiling for Insurance}


\usepackage{hyperref}

\hypersetup{colorlinks,linkcolor={red!50!black},citecolor={blue!50!black},urlcolor={blue!80!black}}

\author{}

\date{}

\begin{document}

\section{Introduction}
\cite{texbook}
\section{Methodology}
The goal of this paper is to answer the following question: "To what extent should insurance companies be allowed to use statistical profiling to determine policy eligibility and pricing?" Since this question does not have a simple answer, subquestions will be answered so that a conclusion to the main question can be drawn. The following subquestions will be used:
\begin{enumerate}
    \item What types of statistical data do insurance companies typically use in risk assessment?
    \item Do any of the statistics used raise concerns about discrimination and if so, who are they discriminant against?
    \item.
    \item.
\end{enumerate}

This risk assessment will be examined through two frameworks: an ethical and a political framework. Within the ethical framework, two main approaches will be used: deontology and consequentialism. 

Deontology is an ethical view that determines the morality of an action based on whether the action itself is moral. "Kant believed moral agents were on duty to comply with what he believed to be a universally valid moral law. According to Kant moral rules are thus supposed to be universal laws. Like the law of gravitational attraction, it is universal and applies to all people equally regardless of status, need or identity" \cite{MSEPA_Syllabus_2025}. This framework is important, because insurance companies make moral choices about fairness in risk assessment. It can be argued that the models that are used unfairly deny insurance based on factors like socioeconomic status, race, gender, etc. 

Consequentialism is an ethical view that states that the only thing that matters is the outcome and not the morality of the action itself \cite{MSEPA_Syllabus_2025}. This is a contradictory view to deontology, as one focuses on the action itself and the other just on the outcome. This view is useful in determining whether the results of statistical profiling are harmful or not. 

According to University of Groningen, "policy is a preferred course of action, with the desired consequence of solving a problem" (2024).



\section{Results}

\bibliographystyle{apacite} % We choose the "plain" reference style
\bibliography{refs} % Entries are in the refs.bib file

\end{document}

